\section{The problem and the proof strategy}
\label{j:introduction}

Let $\mathcal P$ be the probability laws on $\R$ with mean zero,
variance one and finite third absolute moment $\rho(P)=\E_P|X|^3$.
For independent copies with this same law, write $S_n=X_1+\cdots+X_n$
and define
\[
 \Delta_n(P)=\sup_{H\in\R}|\Pp_P(S_n\le H)-\Phi(H/\sqrt n)|,
 \qquad C_n=\sup_{P\in\mathcal P}\frac{\sqrt n\Delta_n(P)}{\rho(P)}.
\]
The threshold $H$ is in the units of the raw sum; $H/\sqrt n$ is its
normal-standardized value. The Berry--Esseen inequality
\cite{Berry1941,Esseen1942,Esseen1945} asserts a universal finite upper
bound on the $C_n$. The sharp-constant problem asks whether this bound
can be taken to be
\[
 \CE=\frac{3+\sqrt{10}}{6\sqrt{2\pi}}=0.409732\ldots.
\]
The lower bound comes from one fixed skew Bernoulli law, rather than
from a sequence of progressively more complicated distributions.
The difficulty is to bound every possible marginal, uniformly in the
sample size.

\begin{theorem}[Sharp iid Berry--Esseen bound]\label{j:main}
For every $P\in\mathcal P$ and every integer $n\ge1$,
\[
 \Delta_n(P)\le\frac{\CE\rho(P)}{\sqrt n}.
\]
The constant is optimal: $\sup_{n\ge1}C_n=\CE$.
Moreover $C_n<\CE$ for every $n\ge47$.
\end{theorem}
The proof combines the analytic arguments below with the finite
inequalities specified by their certificates. Section~\ref{j:composition}
assembles these results and proves strictness using the equality case
of the uniform local exclusion.

\subsection{The three parts of the upper-bound argument}

The proof first establishes $C_j\le\CE$ for $1\le j\le42$.
Suppose a later failure exists and let $n$ be its least index. Then both
deleted predecessors satisfy $C_{n-1},C_{n-2}\le\CE$. A Gaussian
one-step estimate, proved in Section~\ref{j:smoothing:section}, gives
\[
 0<n(C_n-\CE)<\frac{3003}{10250}<0.3,\qquad
 C_n<\frac{73}{175},\qquad n(C_{n-1}-C_n)_+=0.
\]
Two local comparisons exclude precisely such an extremizer:
Proposition~\ref{j:local:finite4346} covers $43\le n\le46$,
and Proposition~\ref{j:local:uniform} is uniform for $n\ge47$.
Thus the finite base and these two comparisons prove the upper bound
at every index. The role of the recurrence is to supply the hypotheses
of the local comparisons, rather than to propagate the sharp constant
by itself.

The mathematical object shared by the comparisons is a nonnegative
support-gap defect. At a fixed index, the optimal ratio is attained.
Its first variation gives a piecewise cubic function which bounds the
conditional probability $y\mapsto\Pp(S_{n-1}\le H-y)$ and agrees
with it on the support. Integrating this contact relation against the
zero-bias density gives a conservation identity. Its nonnegative terms
measure the score above $\CE$, the departure from the extremal
two-point geometry, and the probability entering support gaps. Its
remainder is bounded by Fourier estimates.

Small remainder therefore forces the same extremizing law near two
clusters. Conditioning on the cluster labels gives an exact binomial
count plus centered conditional noise. Couple the conditional sums for
all counts using the same two noise distributions. Replacing one
label-zero summand by a label-one summand increases the total by a
proved positive amount. This orders the threshold crossings and bounds
the entire change in probability by one staircase function of the noise
at a nearest count.

A polynomial of degree four majorizes that staircase on the whole real
line. Its expectation depends only on the first four noise moments;
independence and the acquired support bound those moments explicitly.
The same conditional moments also determine the correction to the third
absolute moment in the Berry--Esseen ratio. Keeping these two calculations
together is the central comparison step.

At the finitely many retained indices, exact binomial gaps and rational
quartic certificates prove the score bound throughout the acquired
parameter regions. For the central large-index range, one explicit
quartic family suffices; four scalar polynomial inequalities give a loss
proportional to every positive noise level, however small. Noise zero is
the actual Bernoulli law and is covered by the binomial bound. Thus the
finite comparison and the uniform tail have different quantitative
reserves, while using the same count coupling and moment argument.

The smallest indices require different quantitative estimates, but
use the same logic: first acquire a valid support and moment domain,
then majorize a probability on that domain. No step assumes that an
unrestricted extremizer has two atoms, or replaces the iid marginal by
an optimizer of a moment relaxation.

\subsection{Where the computer enters}

The reductions leave explicit inequalities on bounded parameter
regions. These include polynomial majorants, enclosures of Fourier
integrals, and inequalities in the probability, noise and phase of a
binomial count. A certificate supplies finite data allowing these
inequalities to be checked on the whole region. Typical data are
rational polynomial coefficients and a finite tree of closed boxes.
The checker proves a valid inequality or an exclusion on each leaf
and verifies that no part of the root domain was lost.

Two mathematical questions precede this computation: why every
possible witness belongs to the specified domain, and why the checked
inequality gives the desired probability bound. Both are proved in
the text. The certificates supply the remaining finite arithmetic.
They are not samples from a class of probability laws. Nor does
numerical precision alone make them rigorous: rational arithmetic and
proved outward enclosures control every rounding error.

Finding a certificate and verifying it are different tasks. For some
majorants, floating-point linear programming proposes coefficients;
for enclosure and count bounds, adaptive subdivision proposes a
finite proof tree. These procedures may fail and need not be trusted.
The acceptance program reconstructs the precise inequality, checks
all of its cells using exact arithmetic, and rejects unresolved cases.
Section~\ref{j:cert:section} explains these constructions and the
common soundness argument. The reproduction package contains the
large coefficient arrays and complete checking programs. Their size
does not enlarge the mathematical assumptions of the proof.

\subsection{Published inputs and sharpness}

For clarity, the external probability estimates used below are stated
here with their full law class.
\begin{proposition}[Affine Berry--Esseen estimates]\label{j:external:affine}
For every $P\in\mathcal P$ and integer $n\ge1$,
\[
 \sqrt n\Delta_n(P)\le
 \min\left\{0.4690\rho,\ 0.3322(\rho+0.429),\
                         0.3031(\rho+0.646)\right\}.
\]
\end{proposition}
These are the iid estimates of \citet{Shevtsova2013}. Standardizing
the summands converts the total-variance normalization there into the
displayed one. Only a finite third absolute moment is assumed. A
putative large normalized score consequently has an explicit upper
bound on $\rho$, which starts several support arguments. All
terminating decimals in this paper denote the corresponding rational
numbers.

\begin{proposition}[Finite binomial comparison and the lower limit]
\label{j:external:binomial}\label{j:external:limit}
For $0<p<1$, let $X=(B-p)/\sqrt{p(1-p)}$ with
$B\sim\operatorname{Bernoulli}(p)$. At every finite $n\ge1$ its
normalized discrepancy is strictly smaller than $\CE$.
For $p_*=(4-\sqrt{10})/2$ and this fixed marginal,
\[
 \lim_{n\to\infty}\frac{\sqrt n\Delta_n(P_*)}{\rho(P_*)}=\CE.
\]
\end{proposition}
The finite assertion is the binomial result of
\citet[Theorem~1]{Schulz2016}; its fixed-law limit is Esseen's lattice
asymptotic \cite{Esseen1956}, in the normalization recorded in
\citet[Remarks~4.1--4.2]{Schulz2016}. Indeed, with $q_*=1-p_*$,
\[
 \rho(P_*)=\frac{p_*^2+q_*^2}{\sqrt{p_*q_*}},\qquad
 \frac{3+q_*-p_*}{6\sqrt{2\pi}(p_*^2+q_*^2)}=\CE.
\]
Thus every coefficient smaller than $\CE$ fails for a sufficiently
large sum of this one law. No interchange of a supremum over laws
with a limit is needed. Together with the upper bound, this proves
that the universal constant is sharp.

The present argument reorganizes the earlier certificate proof
\cite{kolodziejekrodziewicz2026exact}. Its purpose is to give each
analytic reduction once and to expose the precise interfaces between
those reductions and the finite checks. The classical zero-bias
transform \cite{GoldsteinReinert1997} explains the conservation
identity; explicit contact and same-law domain control are needed to
apply it here. The final result concerns the optimal iid constant,
not eventual exact two-point optimality.
